\SourceLeaf{219}%
\noindent Hält man nun \(x\) und \(y\) fest und lässt \(l,\lambda\) variieren, so ergibt die Differentiation
\[
\begin{aligned}
\frac{\mathrm{d}l}{\mathrm{d}\lambda}&=\frac{ab'-ba'}{(a\cos\lambda+a'\sin\lambda)^2+(b\cos\lambda+b'\sin\lambda)^2}\\
&=(ab'-ba')\mathord{.}\frac{\mathrm{d}\sigma^2}{\mathrm{d}s^2}
\end{aligned}
\]
\noindent Je nachdem, ob \(ab'-ba'\) positiv \textit{oder} negativ ist, wachsen \(l\) und \(\lambda\) immer zugleich \textit{oder} ändern sich in entgegengesetztem Sinn. Im ersten Fall sind die Darstellungen 2 und 3 also ähnlich angeordnet, im anderen seitenverkehrt.

Zusammen mit dem zuvor gefundenen Ergebnis folgt daraus, dass die Darstellungen 1 und 3 ähnlich oder seitenverkehrt angeordnet sind, je nachdem, ob \(\frac{ab'-ba'}{h}\) positiv oder negativ ist.

Auf der Fläche mit der Gleichung \(\psi=0\) gilt
\[
e\,\mathrm{d}x+g\,\mathrm{d}y+h\,\mathrm{d}z=0
\]
\noindent und damit auch
\[
(ea+gb+hc)\mathrm{d}t+(ea'+gb'+hc')\mathrm{d}u=0
\]
\noindent unabhängig davon, wie man das Verhältnis von \(\mathrm{d}t\) und \(\mathrm{d}u\) wählt. Daher müssen die folgenden Gleichungen identisch gelten
\[
ea+gb+hc=0,\qquad ea'+gb'+hc'=0
\]
\noindent woraus folgt, dass \(e,g,h\) der Reihe nach zu den Größen \(bc'-cb',ca'-ac',ab'-ba'\) proportional sind. Also gilt
\[
\frac{bc'-cb'}{e}=\frac{ca'-ac'}{g}=\frac{ab'-ba'}{h}
\]
\noindent Man kann daher einen beliebigen dieser drei Ausdrücke heranziehen oder mit der von Natur aus positiven Größe \(ee+gg+hh\) multiplizieren und die daraus entstehende symmetrische Größe
\[
ebc'+gca'+hab'-ecb'-gac'-hba'
\]
\noindent als Kriterium für die ähnliche oder seitenverkehrte Lage der Teile in den Darstellungen 1 und 3 verwenden.

Genauso hängt die ähnliche oder seitenverkehrte Lage der Teile in den Darstellungen 6 und 8 davon ab, ob die Größe
\[
\frac{BC'-CB'}{E}=\frac{CA'-AC'}{G}=\frac{AB'-BA'}{H}
\]
\SourceLeaf{220}%
\noindent oder, wenn man es vorzieht, die symmetrische Größe
\[
EBC'+GCA'+HAB'-ECB'-GAC'-HBA'
\]
\noindent positiv oder negativ ist.

Der Vergleich der Darstellungen 3 und 4 beruht auf ganz ähnlichen Überlegungen wie der von 2 und 3. Die ähnliche oder seitenverkehrte Lage der Teile hängt vom positiven oder negativen Vorzeichen der Größe
\[
(\frac{\mathrm{d}p}{\mathrm{d}t})\mathord{.}(\frac{\mathrm{d}q}{\mathrm{d}u})-(\frac{\mathrm{d}p}{\mathrm{d}u})\mathord{.}(\frac{\mathrm{d}q}{\mathrm{d}t})
\]
\noindent ab. Ebenso bestimmt das positive oder negative Vorzeichen von
\[
(\frac{\mathrm{d}P}{\mathrm{d}T})\mathord{.}(\frac{\mathrm{d}Q}{\mathrm{d}U})-(\frac{\mathrm{d}P}{\mathrm{d}U})\mathord{.}(\frac{\mathrm{d}Q}{\mathrm{d}T})
\]
\noindent die ähnliche oder seitenverkehrte Lage der Teile in den Darstellungen 5 und 6.

Beim Vergleich der Darstellungen 4 und 5 können wir schließlich auf die Untersuchung in Artikel 8 zurückgreifen. Daraus geht hervor, dass ihre kleinsten Teile ähnlich oder seitenverkehrt angeordnet sind, je nachdem, ob man die erste oder die zweite Lösung wählt, also entweder
\[
P+iQ=f(p+iq)\quad\text{und}\quad P-iQ=f'(p-iq)
\]
\Needspace{5\baselineskip}
\noindent oder
\[
P+iQ=f(p-iq)\quad\text{und}\quad P-iQ=f'(p+iq)
\]
\noindent setzt.

Wir folgern daraus, dass für eine Abbildung auf der Fläche mit der Gleichung \(\Psi=0\), deren kleinste Teile dem Urbild auf der Fläche mit der Gleichung \(\psi=0\) ähnlich und auch ähnlich angeordnet sein sollen, die Anzahl der negativen unter den folgenden vier Größen maßgeblich ist
\[
\begin{gathered}
\frac{ab'-ba'}{h},\quad (\frac{\mathrm{d}p}{\mathrm{d}t})\mathord{.}(\frac{\mathrm{d}q}{\mathrm{d}u})-(\frac{\mathrm{d}p}{\mathrm{d}u})\mathord{.}(\frac{\mathrm{d}q}{\mathrm{d}t}),\\
(\frac{\mathrm{d}P}{\mathrm{d}T})\mathord{.}(\frac{\mathrm{d}Q}{\mathrm{d}U})-(\frac{\mathrm{d}P}{\mathrm{d}U})\mathord{.}(\frac{\mathrm{d}Q}{\mathrm{d}T}),\quad \frac{AB'-BA'}{H}
\end{gathered}
\]
\noindent Ist keine oder eine gerade Anzahl von ihnen negativ, muss man die erste Lösung wählen. Sind eine oder drei negativ, muss man die zweite wählen. Die jeweils entgegengesetzte Wahl ergibt stets eine seitenverkehrte Ähnlichkeit.
\par\indent\SourceLeaf{221}%
Bezeichnet man die vier obigen Größen der Reihe nach mit \(r,s,S,R\), so lässt sich außerdem zeigen, dass immer
\[
\frac{r\BareRoot{(ee+gg+hh)}}{s}=\pm n,\qquad\frac{R\BareRoot{(EE+GG+HH)}}{S}=\pm N
\]
\noindent gilt, wobei \(n\) und \(N\) die Bedeutung aus Artikel 5 haben. Wir lassen den leicht zu findenden Beweis dieses Satzes hier beiseite, weil wir ihn für unseren Zweck nicht weiter brauchen.

